\subsection{Global Navigation Satellite System (GNSS)}

Os sistemas globais de navegação por satélite (GNSS) realizam a determinação da posição, velocidade e tempo para diversas aplicações, tanto civis quanto militares.
Esses sistemas incluem constelações como o \emph{GPS} (EUA), \emph{GLONASS} (Rússia), \emph{Galileo} (Europa), \emph{BeiDou} (China), \emph{NavIC} (Índia) e \emph{QZSS} (Japão), que operam em diferentes configurações orbitais para oferecer cobertura global e regional.
A precisão dos GNSS depende da sincronização de sinais entre satélites e receptores terrestres.

A disponibilidade e desses sistemas podem ser comprometidas por fatores como obstruções no sinal, interferências atmosféricas e limitações impostas por receptores de antena única.
Para mitigar essas limitações, a integração do GNSS com sistemas inerciais de navegação (INS) tem sido estudada, resultando no desenvolvimento de abordagens que utilizam filtros de Kalman para estimar erros e aprimorar a precisão da navegação, mesmo em ambientes onde os sinais GNSS são intermitentes ou ausentes. 
Com a evolução dos sistemas SATNAV, avanços foram trazidos na transmissão de sinais em múltiplas frequências e no uso de correções baseadas em redes de estações terrestres, permitindo confiabilidade dos sistemas de navegação global \cite{kaplan2017}.

% \subsubsection{Integrações loosely, tightly e variantes}

% Os sistemas de navegação integrados GNSS evoluíram para atender às exigências de precisão e robustez em aplicações terrestres, aéreas e espaciais. A integração \textit{loosely coupled}, representa a forma mais simples de fusão de dados entre GNSS e INS, onde as estimativas de posição e velocidade do GNSS são utilizadas para corrigir as medições do sistema inercial por meio de um filtro de Kalman. Esse método apresenta limitações quando há perda temporária dos sinais GNSS, como em áreas urbanas com obstruções ou em ambientes subterrâneos, pois depende diretamente da disponibilidade de múltiplos satélites para atualização contínua \cite{falco2017}.

% Como explica \cite{falco2017}, para superar essas restrições, a integração \textit{tightly coupled} combina diretamente os dados brutos do GNSS, como pseudodistâncias e medições Doppler, com as leituras do INS, possibilitando que o sistema continue a operar mesmo quando menos de quatro satélites estão disponíveis. Essa técnica reduz erros da navegação inercial e melhora a estabilidade da solução em condições adversas. 

% Para aplicações que exigem ainda mais precisão e resistência a falhas de sinal, o \textit{deeply coupled} ou \textit{ultra-tight coupling} leva a fusão dos sistemas a um nível mais avançado, onde o próprio processamento do sinal GNSS é ajustado com base nos dados do INS, permitindo um rastreamento mais robusto dos satélites em situações de bloqueio parcial dos sinais, como em operações militares, navegação autônoma e missões espaciais. Essa abordagem requer maior capacidade computacional e algoritmos mais sofisticados, mas proporciona maior resiliência à interferência e degradação dos sinais GNSS, tornando-se essencial para veículos autônomos e sistemas de posicionamento crítico \cite{falco2017}. 

\subsection{Pseudolites}

Os pseudolites são transmissores terrestres que simulam sinais de satélites GNSS, permitindo que os receptores GPS utilizem esses sinais para posicionamento mesmo em ambientes onde os sinais dos satélites são fracos, bloqueados por obstáculos ou inexistentes \cite{kim2014}.
O principal desafio na implementação de sistemas baseados em pseudolites está na compatibilidade com receptores GPS antigos e que normalmente não são projetados para interpretar sinais provenientes de fontes terrestres.
Para contornar essa limitação, é possível modificar os sinais emitidos pelos pseudolites para que imitem os dados de posição e o movimento dos satélites ao longo do tempo, ou seja, satélites simulados.

Com isso, os receptores GPS processem os sinais sem necessidade de ajustes no \emph{hardware} ou \emph{firmware}.
No entanto, essa solução apresenta desafios, pois a posição calculada inicialmente pelo receptor GPS tende a ser incorreta devido a atrasos adicionais introduzidos na propagação do sinal entre cada pseudolite e o receptor, resultando em um processo de pós-cálculo é necessário para corrigir a posição estimada, utilizando informações sobre os sinais simulados, os pseudolites e o próprio processo de posicionamento do GPS.

% A sincronização dos pseudolites deve ser bem calculada para garantir a precisão do sistema, sendo realizada por meio de uma estação de referência que coleta medições dos sinais transmitidos e gera entradas de controle para corrigir os relógios dos pseudolites. Além disso, um servidor de dados é necessário para fornecer de forma contínua as informações sobre a posição dos pseudolites e os parâmetros necessários para o receptor interpretar corretamente os sinais recebidos. Esse servidor também se comunica com os pseudolites para garantir que a transmissão dos sinais ocorra de maneira sincronizada e coerente com o sistema GPS.

% \subsection{Arquitetura do pseudolite}

% A arquitetura do sistema de pseudolites é composta por quatro elementos principais: os próprios pseudolites, a estação de referência, o servidor de dados e o terminal do usuário. Os pseudolites geram sinais GPS simulados na banda L1, cuja frequência é de 1.575 $\pm$ 5 MHz, conforme explica \cite{wang2018} e transmitem esses sinais aos receptores convencionais, utilizando os dados fornecidos pelo servidor.
% A estação de referência é responsável por manter a sincronização entre os pseudolites, garantindo que as diferenças de tempo entre os sinais não comprometam a precisão do posicionamento. O servidor de dados armazena informações sobre a configuração do sistema e distribui os dados necessários para a operação dos pseudolites e para a correção da posição estimada pelo receptor GPS. O terminal do usuário, que pode ser um dispositivo convencional com capacidade de recepção GPS, recebe os sinais dos pseudolites e processa os dados em conjunto com as informações fornecidas pelo servidor para obter uma posição corrigida.
% Para melhorar a precisão do posicionamento, um mecanismo de correção de erros nos sinais dos pseudolites pode ser implementado por meio da aplicação de filtros digitais e algoritmos de ajuste de tempo, que compensam atrasos e distorções causados pelo ambiente \cite{kim2014}.